% =====================================================================
%  INTRODUÇÃO
%  Importe com % =====================================================================
%  INTRODUÇÃO
%  Importe com % =====================================================================
%  INTRODUÇÃO
%  Importe com % =====================================================================
%  INTRODUÇÃO
%  Importe com \input{introducao} a partir do seu .tex principal.
%  As chaves \cite{} seguem o arquivo referencias.bib deste trabalho.
% =====================================================================

\chapter{Introdução}
\label{ch:introducao}
As redes de cápsulas prometem equivariância. A propriedade vem da forma como a cápsula representa uma entidade: a orientação do vetor codifica a pose. Com a pose explícita, a rede aprende a relação entre as características detectadas nas cápsulas primárias. Dessa relação viria a resistência a rotações e a outras transformações ausentes do treinamento.

A promessa é de 2017. Desde então, augmentação e regularização deram aos modelos convolucionais e de atenção boa parte da tolerância que se atribuía à arquitetura. Fica a dúvida se a vantagem das cápsulas ainda existe.

Para esclarecer essa dúvida, este trabalho compara três arquiteturas: Efficient-CapsNet, ResNet-18 e DeiT-Tiny. Treinamos as três do zero em CIFAR-10 e D-Fire, sob três regimes de augmentação (nenhuma, padrão, forte), três frações de dados (33\%, 66\%, 100\%) e três sementes.

A comparação só é justa se nada além da arquitetura variar. As três usam os mesmos hiperparâmetros, as mesmas sementes e as mesmas partições. Comparamos sob dois orçamentos: mesmo número de épocas e mesmo tempo de treino.

O teste da equivariância vem depois do treino. Avaliamos cada modelo sob transformações que ele nunca viu e medimos quanto da acurácia limpa cada arquitetura retém.
 a partir do seu .tex principal.
%  As chaves \cite{} seguem o arquivo referencias.bib deste trabalho.
% =====================================================================

\chapter{Introdução}
\label{ch:introducao}
As redes de cápsulas prometem equivariância. A propriedade vem da forma como a cápsula representa uma entidade: a orientação do vetor codifica a pose. Com a pose explícita, a rede aprende a relação entre as características detectadas nas cápsulas primárias. Dessa relação viria a resistência a rotações e a outras transformações ausentes do treinamento.

A promessa é de 2017. Desde então, augmentação e regularização deram aos modelos convolucionais e de atenção boa parte da tolerância que se atribuía à arquitetura. Fica a dúvida se a vantagem das cápsulas ainda existe.

Para esclarecer essa dúvida, este trabalho compara três arquiteturas: Efficient-CapsNet, ResNet-18 e DeiT-Tiny. Treinamos as três do zero em CIFAR-10 e D-Fire, sob três regimes de augmentação (nenhuma, padrão, forte), três frações de dados (33\%, 66\%, 100\%) e três sementes.

A comparação só é justa se nada além da arquitetura variar. As três usam os mesmos hiperparâmetros, as mesmas sementes e as mesmas partições. Comparamos sob dois orçamentos: mesmo número de épocas e mesmo tempo de treino.

O teste da equivariância vem depois do treino. Avaliamos cada modelo sob transformações que ele nunca viu e medimos quanto da acurácia limpa cada arquitetura retém.
 a partir do seu .tex principal.
%  As chaves \cite{} seguem o arquivo referencias.bib deste trabalho.
% =====================================================================

\chapter{Introdução}
\label{ch:introducao}
As redes de cápsulas prometem equivariância. A propriedade vem da forma como a cápsula representa uma entidade: a orientação do vetor codifica a pose. Com a pose explícita, a rede aprende a relação entre as características detectadas nas cápsulas primárias. Dessa relação viria a resistência a rotações e a outras transformações ausentes do treinamento.

A promessa é de 2017. Desde então, augmentação e regularização deram aos modelos convolucionais e de atenção boa parte da tolerância que se atribuía à arquitetura. Fica a dúvida se a vantagem das cápsulas ainda existe.

Para esclarecer essa dúvida, este trabalho compara três arquiteturas: Efficient-CapsNet, ResNet-18 e DeiT-Tiny. Treinamos as três do zero em CIFAR-10 e D-Fire, sob três regimes de augmentação (nenhuma, padrão, forte), três frações de dados (33\%, 66\%, 100\%) e três sementes.

A comparação só é justa se nada além da arquitetura variar. As três usam os mesmos hiperparâmetros, as mesmas sementes e as mesmas partições. Comparamos sob dois orçamentos: mesmo número de épocas e mesmo tempo de treino.

O teste da equivariância vem depois do treino. Avaliamos cada modelo sob transformações que ele nunca viu e medimos quanto da acurácia limpa cada arquitetura retém.
 a partir do seu .tex principal.
%  As chaves \cite{} seguem o arquivo referencias.bib deste trabalho.
% =====================================================================

\chapter{Introdução}
\label{ch:introducao}
As redes de cápsulas prometem equivariância. A propriedade vem da forma como a cápsula representa uma entidade: a orientação do vetor codifica a pose. Com a pose explícita, a rede aprende a relação entre as características detectadas nas cápsulas primárias. Dessa relação viria a resistência a rotações e a outras transformações ausentes do treinamento.

A promessa é de 2017. Desde então, augmentação e regularização deram aos modelos convolucionais e de atenção boa parte da tolerância que se atribuía à arquitetura. Fica a dúvida se a vantagem das cápsulas ainda existe.

Para esclarecer essa dúvida, este trabalho compara três arquiteturas: Efficient-CapsNet, ResNet-18 e DeiT-Tiny. Treinamos as três do zero em CIFAR-10 e D-Fire, sob três regimes de augmentação (nenhuma, padrão, forte), três frações de dados (33\%, 66\%, 100\%) e três sementes.

A comparação só é justa se nada além da arquitetura variar. As três usam os mesmos hiperparâmetros, as mesmas sementes e as mesmas partições. Comparamos sob dois orçamentos: mesmo número de épocas e mesmo tempo de treino.

O teste da equivariância vem depois do treino. Avaliamos cada modelo sob transformações que ele nunca viu e medimos quanto da acurácia limpa cada arquitetura retém.
